% TOP_PROBLEM: {"id": 30006883, "title": "The Word Problem for One-Relation Monoids", "category_id": 4, "category": {"id": 4, "name": "algebra", "display_name": "Algebra", "description": "Group theory, ring theory, field theory, and algebraic structures.", "slug": "algebra", "order_index": 4, "created_at": "2026-07-31T15:26:25.671Z"}, "status": "open"}
% FIRST_POSTED: 2026-09-27
% !TeX program = lualatex
% ATTEMPT_STATUS: unresolved
\documentclass[11pt]{article}
\usepackage[margin=25mm]{geometry}
\usepackage{fontspec}
\setmainfont{Latin Modern Roman}
\usepackage{amsmath,amssymb,mathtools}
\usepackage{unicode-math}
\setmathfont{Latin Modern Math}
\usepackage{xurl,hyperref}
\hypersetup{colorlinks=true,urlcolor=blue}
\setcounter{secnumdepth}{0}
\setlength{\parindent}{0pt}
\setlength{\parskip}{5pt}
\setlength{\emergencystretch}{3em}
\begin{document}
\section*{\texorpdfstring{The Word Problem for One-Relation Monoids}{The Word Problem for One-Relation Monoids}}\label{30006883}
\subsection{Definitions and mathematical statement}
For a finite alphabet $A=\{a_1,\ldots,a_n\}$, the free monoid $A^*$ consists of finite words, including the empty word, with concatenation. A one-relation monoid is
\[
 M=\langle A\mid u=v\rangle=A^*/{\sim},
\]
where $\sim$ is the least equivalence relation compatible with concatenation and containing $u\sim v$. Its word problem is to decide, given $w,w'$, whether $w\sim w'$. The catalogue combines two formulations: decidability for every fixed one-relation presentation, and an algorithm taking the presentation itself as input. The latter is a uniform decision problem and is logically stronger than an unqualified assertion that each fixed presentation has some algorithm. Both interpretations should be kept distinct; an algorithm announced for the general case must specify which quantifiers it establishes.
\par\smallskip\textit{Formulation and concept references:} \href{https://arxiv.org/abs/2105.02853}{[S1]}.

\subsection{Short English statement}
Can equality of words always be decided in monoids defined by one relation, and can this be done uniformly from the presentation?
\subsection{Research attempt}
\paragraph{Scope and literature audit (27 September 2026).}
Equality means a finite sequence of substitutions inside arbitrary left
and right contexts, using either direction of the single relation.
The empty word is allowed. Cancellation, inversion and changes of
presentation must be justified; they are not free operations on words
in a monoid. The uniform version asks for one algorithm taking
$(A,u,v,w,w')$, not a different unspecified algorithm for each pair
$(u,v)$. Every positive algorithm below is uniform on its stated class.

Nyberg-Brodda's primary survey [S1] explicitly treats the general problem
as open and warns about inaccurate reductions in the literature. Its
accepted version was inspected. Recent papers on Diophantine problems,
submonoid membership and one-relator \emph{inverse} monoids address
other decision problems; their titles do not supply an undecidable
ordinary word problem for a one-relation monoid. The present attempt
uses none of those claims as a general solution.

The attack is to orient the relation and identify exactly when terminating
reduction gives an equality test. This yields several complete classes,
an effective sufficient confluence test, and explicit counterchecks of
the two common shortcuts: canceling letters and choosing an arbitrary
irreducible representative.

\paragraph{The positive instances are uniformly enumerable.}
Put an edge between words $p u q$ and $p v q$, for all $p,q\in A^*$.
The connected components of this undirected graph are exactly the
congruence classes. To see this, graph connectivity is an equivalence
relation containing $u=v$ and is preserved by adding contexts. Conversely
every edge belongs to the defining congruence. The two minimality
statements establish equality of the relations.

Each word has finitely many neighbors, effectively listable by scanning
its possible occurrence positions. This remains true if a relation side
is empty: a word of length $L$ has $L+1$ insertion positions. Breadth-first
search starting from $w$ therefore eventually finds $w'$ if the two
are congruent. It also returns a finite checkable substitution certificate.
It need not halt when they are unequal. Running searches from both ends
does not change that issue; two disjoint infinite components can each
continue producing new vertices forever.

If $u=v$ as literal strings, there are no nontrivial substitutions and
free-word equality is the answer. Assume henceforth that the two sides
differ. Alphabet symbols not occurring in the relation are still retained
in every algorithm, rather than silently removed from the input.

\paragraph{A complete uniform algorithm for length-preserving relations.}
\textbf{Proposition 359.1.} If $|u|=|v|$, the word problem is uniformly
decidable from the presentation.

\emph{Proof.} Substitutions preserve length. Unequal input lengths
therefore prove inequality immediately. If both lengths are $L$, the
entire component of either input is contained in the finite set $A^L$.
Enumerate this graph and run breadth-first search, stopping when its
finite reachable set is exhausted. Membership of $w'$ in that set is
exactly congruence, by the preceding graph characterization.
\hfill$\square$

For $a=|A|\ge2$ there are at most $a^L$ vertices. A direct implementation
scans $O(L)$ possible positions per side, compares words and stores
visited strings of length $L$. Thus a deliberately coarse bound is
$a^L$ times a polynomial in $L+|u|+a$, both in time and in total stored
bits. This is a decision theorem, not a polynomial-time claim. For a
one-letter alphabet there is only one word of each length. The empty
alphabet has just its empty word and is immediate as well.

The proof depends on all paths preserving length. For unequal relation
lengths, the existence of short endpoints does not bound intermediate
lengths. Imposing an arbitrary length cutoff then changes a complete
search into an incomplete one.

\paragraph{Orienting one rule gives termination, but not automatically a normal form.}
Fix an order on $A$ and use shortlex order on words: compare length
first, then lexicographic order among words of equal length. Orient the
larger relation side $l$ toward the smaller side $r$. A replacement
$p l q\to p r q$ strictly decreases shortlex order. If lengths differ,
the total length decreases; if they agree, the first changed letter
lies inside the occurrence and decreases lexicographically.

All descendants of a word of length $L$ lie in the finite set of words
of length at most $L$. Consequently every forward reduction terminates.
When $|l|>|r|$, there are at most $L/(|l|-|r|)$ steps; a naive rescan
implementation is polynomial in the explicit input length. In the
length-preserving case a strictly decreasing sequence can have
exponentially many terms, and no polynomial bound follows merely
from this order.

A terminal word contains no occurrence of $l$ and is called irreducible.
An arbitrary deterministic choice of occurrence computes \emph{an}
irreducible word. To make it an equality invariant requires confluence:
any two forward descendants of a common ancestor must have a common
forward descendant. Termination and confluence together give a unique
irreducible word, which we denote by $\operatorname{nf}(w)$.

\paragraph{Why confluence makes reduction a complete decision procedure.}
Suppose the oriented system is terminating and confluent. Every word
has an irreducible descendant by termination; two such descendants must
coincide by confluence. If $x\to y$ is one elementary reduction, then
$x$ and $y$ have the same normal form. The same is true of its inverse.
Along every finite congruence path the normal form is therefore unchanged.
Conversely, reducing two inputs to a common word gives a congruence path
between them by reversing one of the reductions. Hence
\[
 w\sim w'\quad\Longleftrightarrow\quad
 \operatorname{nf}(w)=\operatorname{nf}(w').
 \tag{359.1}
\]
This proof covers substitutions in either direction even though the
algorithm itself only uses the decreasing direction.

Here is the termination-to-confluence lemma in the form needed below.
A terminating relation is confluent if every one-step peak
$x\leftarrow z\to y$ has joinable endpoints. Prove this by well-founded
induction on $z$. Two reduction paths either start with the same step,
in which case the induction hypothesis below that step joins them, or
start with a peak. Join its endpoints at $v$. Induct below each endpoint
to join its given descendant with $v$; then induct below $v$ to join
the resulting descendants. All applications are strictly below $z$.
This is the usual proof of Newman's lemma; its decreasing hypothesis
is available here, rather than assumed for an arbitrary completion.

\paragraph{A complete class: the larger side has no proper self-overlap.}
A nonempty word $l$ is unbordered if no nonempty proper prefix is also
a suffix. This is decidable by comparing its at most $|l|-1$ possible
border lengths.

\textbf{Proposition 359.2.} If the shortlex-larger side $l$ is unbordered,
the single-rule system $l\to r$ is confluent and uniformly decides
the presented word problem by (359.1).

\emph{Proof.} Consider two occurrences of $l$ in a word. If their
intervals coincide, the two reductions agree. If they overlap properly,
the overlap is a nonempty proper suffix and prefix of $l$, contradicting
unborderedness. Otherwise they are disjoint. Replacing either occurrence
leaves a corresponding occurrence of the other, and doing both replacements
in either order gives the same word. New occurrences created by a
replacement are irrelevant: the original other occurrence still gives
this specific joining path. Thus all one-step peaks join, and termination
plus the preceding lemma gives confluence.\hfill$\square$

For example $ab\to1$ has an unbordered left side. So does any relation
whose larger side has all letters distinct. The hypothesis is sufficient,
not necessary: $aa\to a$ has an overlapping left side but plainly still
has a unique irreducible word on each block of $a$'s.

\paragraph{All remaining critical peaks can be checked effectively.}
For a single rule, disjoint occurrences have already been handled.
A proper overlap is described by
\[
 l=ab=bc,\qquad 0<|b|<|l|,
\]
where the word $abc$ contains the two overlapping occurrences of $l$.
The two immediate descendants are $rc$ and $ar$. There are at most
$|l|-1$ such overlaps. There are no proper containment peaks between
two equal-length copies of the same left side.

The descendants of each endpoint lie in a finite, explicitly bounded
set of words, because our fixed rule is nonincreasing in length and
terminating. Enumerate their forward descendants and test whether the
two sets intersect. If every overlap pair joins, the same joins work
inside arbitrary surrounding contexts. All one-step peaks then join,
so the preceding proof establishes confluence. Conversely confluence
requires each of these finitely many pairs to join.

This gives a terminating algorithm to decide whether this particular
one-rule orientation is confluent. If it is, it also supplies a uniform
word-problem algorithm. A negative answer says only that this orientation
has inadequate normal forms. It does not imply undecidability of the
monoid, or rule out another finite or infinite effective rewriting system.

\paragraph{An exact counterexample to the irreducible-word shortcut.}
Take the relation $aba=ab$ and orient it as $aba\to ab$. It is strictly
length reducing. The word $ababa$ has two overlapping occurrences:
\[
 ababa\longrightarrow abba,
 \qquad
 ababa\longrightarrow abab\longrightarrow abb.
\]
Both $abba$ and $abb$ contain no $aba$, so both are irreducible. They
are different words but are congruent through the displayed ancestor.
Thus comparing arbitrary irreducible reductions can report a false
negative, even with only one defining relation and a strictly decreasing
length measure.

The finite critical-peak test detects precisely this failure. Adding
an extra rule to join these two outputs is a possible next step, but
that creates further overlaps requiring their own analysis. A promise
that completion eventually stabilizes would be an additional theorem,
not a consequence of having started with one relation.

\paragraph{A worked noncancellative example: the bicyclic monoid.}
For $M=\langle a,b\mid ab=1\rangle$, Proposition 359.2 proves that
the unique irreducible words are
\[
 b^i a^j,\qquad i,j\ge0.
\]
Indeed, a word with no adjacent $ab$ has all its $b$'s before all its
$a$'s: if some $a$ preceded some later $b$, the first transition from
an $a$-block to a $b$-block would contain $ab$. Conversely the displayed
words contain no reducible factor. Their multiplication is explicit:
\[
 (b^ia^j)(b^ka^l)=b^{i+k-t}a^{j+l-t},
 \qquad t=\min(j,k).
 \tag{359.2}
\]
Only the central $a^jb^k$ cancels, exactly $t$ times.

This example refutes a reduction to the one-relator \emph{group}
with the same relation. That group also has $ba=1$, but in the monoid
$ba$ and the empty word are distinct irreducibles. An independent
representation makes the failure concrete. On $\mathbb N$ let
$a(n)=\max(n-1,0)$ and $b(n)=n+1$, composing maps right to left.
Then $ab$ is the identity, whereas $ba(0)=1$. Hence the relation holds
in this transformation monoid but $ba=1$ fails. No group-completion
argument can preserve all monoid equalities and inequalities here.

Nor can one cancel the outer $a$ from $aba=a$, although that equality
is valid because $ab=1$. Left cancellation would incorrectly give
$ba=1$. This is a concrete reason to check cancellativity before using
cancellation in a proposed algorithm.

\paragraph{The complete one-generator case, including an empty side.}
Let the alphabet be $\{a\}$ and the relation be $a^r=a^s$, with
$r>s\ge0$. Put $d=r-s>0$. Reducing $a^r$ to $a^s$ subtracts $d$
from every length at least $r$. Define
\[
 f(L)=\begin{cases}
 L,&L<s,\\
 s+((L-s)\bmod d),&L\ge s.
 \end{cases}
\]
Repeated reduction ends at $a^{f(L)}$: for $s\le L<r$ this formula
already equals $L$, and for $L\ge r$ each step subtracts $d$ until
that interval is reached. Every occurrence replacement has the same
length effect, so every reduction path has this endpoint.

A relation step in either direction preserves $f$. Indeed its two
lengths are $r+h$ and $s+h$ for some $h\ge0$, both at least $s$ and
congruent modulo $d$. Thus equality is equivalent to equality of these
normal lengths. This proves a uniform decision algorithm, including
$s=0$, when the monoid is cyclic of order $r$. If $r=s$, literal
length equality is the answer instead. The arithmetic formula is not
an assertion that arbitrary multi-letter words are determined by their
lengths.

\paragraph{Useful negative certificates and their limit.}
Every assignment of weights $c_a\in\mathbb Z$ to letters for which
$\sum_a c_a|u|_a=\sum_a c_a|v|_a$ gives a monoid homomorphism to
$(\mathbb Z,+)$. Different total weights certify inequality. The same
is true of any explicitly supplied finite monoid and any letter
assignment satisfying the relation: evaluate both input words there.
These are sound tests, and finite monoids and their multiplication
tables can be exhaustively enumerated and checked for associativity.

Running all such finite tests in parallel with congruence-path search
would decide the problem for a presentation known to be residually
finite. That hypothesis cannot be dropped. The bicyclic monoid already
shows why: in every finite monoid $ab=1$ implies $ba=1$. To prove this,
left multiplication by $a$ is surjective because $a(bx)=x$, hence
injective on the finite set. Since $a(ba)=a=a1$, injectivity gives
$ba=1$. Thus no finite-monoid homomorphism separates its distinct
words $ba$ and $1$. Its word problem is nevertheless decidable by the
normal forms above. Finite-quotient separation is a sufficient strategy,
not an automatic property or a necessary route to decidability.

\paragraph{Diagnostics and first unresolved step.}
The companion exact checks enumerate overlap descendants, recover the
two irreducibles for $aba=ab$, compare the bicyclic multiplication formula
with actual reduction, and check the one-letter normal-length invariant.
All calculations concern finite words and exact substitutions; no finite
cutoff is interpreted as proving inequality in an unrestricted search.

\textbf{Outcome: UNRESOLVED.} The attempt supplies uniform decision
algorithms for length-preserving relations, unbordered decreasing sides,
all presentations passing the finite one-rule confluence test, and the
one-generator case. In the general unequal-length case the first gap
is a complete effective way to control nonconfluent reductions or all
possible expansions. Neither terminating one-way reduction, a finite
quotient search, nor passing to a one-relator group supplies that step.

\subsection{Sources}
\begin{itemize}
\item[S1]The Word Problem for One-relation Monoids: A Survey (Carl-Fredrik Nyberg-Brodda).
\url{https://arxiv.org/abs/2105.02853}.
\textit{Formulation source.}
\item[E1]Carl-Fredrik Nyberg-Brodda, \emph{Undecidability of the
Diophantine problem for one-relator groups and one-relation monoids},
arXiv:2608.01983v1, 3 August 2026.
\url{https://arxiv.org/abs/2608.01983}.
\textit{Primary abstract inspected 27 September 2026. Solving equations
with unknowns is distinguished from equality of two specified words.}
\item[E2]Robert D. Gray and Catherine Reilly, \emph{The word problem
for two-generator one-relator inverse monoids}, arXiv:2608.04650v1,
5 August 2026.
\url{https://arxiv.org/abs/2608.04650}.
\textit{Primary abstract inspected 27 September 2026. The inverse-monoid
category and prefix-membership problem are not the ordinary monoid
word problem of the catalogue.}
\end{itemize}
\end{document}
